\chapter{Теория}

\section{netstat Windows}

\begin{figure}[H] 
    \centering
    \includegraphics[width=0.75\textwidth]{ns.png}
    \caption{Результат работы netstat}
\end{figure}

Результат вывода команды содержит следующие элементы:
\begin{enumerate}[label=\arabic*.]
  \item \textbf{Протокол активного сокета}, TCP в нашем случае.
  \item \textbf{Адрес сокета} (имя\_хоста:порт, либо сеть:порт).
  \item \textbf{Удалённый адрес и порт}, с которым установлено соединение. В данном случае \texttt{https} означает порт \texttt{443} (стандартный порт для HTTPS).
  \item \textbf{Состояние соединения}. Разберём самые часто встречающиеся виды соединений:
  \begin{itemize}[label=--]
    \item \textbf{ESTABLISHED} --- соединение установлено и данные могут передаваться между двумя сторонами.
    \item \textbf{TIME\_WAIT} --- ожидание определённого времени, чтобы убедиться, что удалённая сторона получила подтверждение о прекращении соединения.
    \item \textbf{LISTEN} --- ожидание входящих соединений.
    \item \textbf{CLOSED} --- на данный момент порт не используется каким-либо приложением, но может быть выдан ему в любой момент.
    \item \textbf{SYN\_SENT} --- активно пытается установить соединение.
    \item \textbf{SYN\_RECEIVED} --- идёт начальная синхронизация соединения.
  \end{itemize}
\end{enumerate}

\subsection{Ключи}

\subsubsection{-a}

Показывает состояние всех сокетов. Обычно сокеты, используемые серверными
процессами, не отображаются.

\subsubsection{-n}
Вывод активных подключений TCP с отображением адресов и номеров портов в
числовом формате без попыток определения имен.

\subsubsection{-o}
Вывод активных подключений TCP и включение кода процесса (PID) для каждого
подключения.

\subsubsection{-r}
Показывает таблицы маршрутизации.

\subsubsection{-p}
Вывод подключений для указанного протокола (tcp).



\endinput